% Slug: ch33-heap-graph
% Source: chapters/ch33.tex:29 (tikzpicture env start line)
% Original label: \label{fig:ch33-heap-graph}
% Original caption (LaTeX source):
%   Two roots reach A, B, C, and D through shared and direct edges, while E and F have no path from any root and are garbage.
% Original caption (plain text):
%   Two roots reach A, B, C, and D through shared and direct edges, while E and F have no path from any root and are garbage.
% Required TikZ setup: \usepackage{tikz} + \usetikzlibrary{arrows.meta, positioning, shapes.geometric, fit} (loaded by support/compiler-pseudocode.sty).
% Shared styles used in this picture: `basic block', `cfg edge'.
%   Defined in support/compiler-pseudocode.sty: `basic block', `cfg edge', `back edge', `block label'.
% Usage: % Slug: ch33-heap-graph
% Source: chapters/ch33.tex:29 (tikzpicture env start line)
% Original label: \label{fig:ch33-heap-graph}
% Original caption (LaTeX source):
%   Two roots reach A, B, C, and D through shared and direct edges, while E and F have no path from any root and are garbage.
% Original caption (plain text):
%   Two roots reach A, B, C, and D through shared and direct edges, while E and F have no path from any root and are garbage.
% Required TikZ setup: \usepackage{tikz} + \usetikzlibrary{arrows.meta, positioning, shapes.geometric, fit} (loaded by support/compiler-pseudocode.sty).
% Shared styles used in this picture: `basic block', `cfg edge'.
%   Defined in support/compiler-pseudocode.sty: `basic block', `cfg edge', `back edge', `block label'.
% Usage: % Slug: ch33-heap-graph
% Source: chapters/ch33.tex:29 (tikzpicture env start line)
% Original label: \label{fig:ch33-heap-graph}
% Original caption (LaTeX source):
%   Two roots reach A, B, C, and D through shared and direct edges, while E and F have no path from any root and are garbage.
% Original caption (plain text):
%   Two roots reach A, B, C, and D through shared and direct edges, while E and F have no path from any root and are garbage.
% Required TikZ setup: \usepackage{tikz} + \usetikzlibrary{arrows.meta, positioning, shapes.geometric, fit} (loaded by support/compiler-pseudocode.sty).
% Shared styles used in this picture: `basic block', `cfg edge'.
%   Defined in support/compiler-pseudocode.sty: `basic block', `cfg edge', `back edge', `block label'.
% Usage: % Slug: ch33-heap-graph
% Source: chapters/ch33.tex:29 (tikzpicture env start line)
% Original label: \label{fig:ch33-heap-graph}
% Original caption (LaTeX source):
%   Two roots reach A, B, C, and D through shared and direct edges, while E and F have no path from any root and are garbage.
% Original caption (plain text):
%   Two roots reach A, B, C, and D through shared and direct edges, while E and F have no path from any root and are garbage.
% Required TikZ setup: \usepackage{tikz} + \usetikzlibrary{arrows.meta, positioning, shapes.geometric, fit} (loaded by support/compiler-pseudocode.sty).
% Shared styles used in this picture: `basic block', `cfg edge'.
%   Defined in support/compiler-pseudocode.sty: `basic block', `cfg edge', `back edge', `block label'.
% Usage: \input{graphs-code/ch33-heap-graph.tikz.tex} inside a figure environment
%   (wrap with \begin{figure} ... \caption{...} \label{fig:ch33-heap-graph} ... \end{figure}).
% Note: tikzpicture body below is verbatim from the source; do not reformat.
\begin{tikzpicture}[node distance=1.0cm and 1.5cm]
\node[basic block] (r1) {root R1};
\node[basic block, below=of r1] (a) {A};
\node[basic block, below left=of a] (b) {B};
\node[basic block, below right=of a] (c) {C};
\node[basic block, right=of r1] (r2) {root R2};
\node[basic block, below=of r2] (d) {D};
\node[basic block, below=of d] (e) {E dead};
\node[basic block, left=of e] (f) {F dead};
\draw[cfg edge] (r1) -- (a);
\draw[cfg edge] (a) -- (b);
\draw[cfg edge] (a) -- (c);
\draw[cfg edge] (r2) -- (d);
\draw[cfg edge] (d) -- (a);
\end{tikzpicture}
 inside a figure environment
%   (wrap with \begin{figure} ... \caption{...} \label{fig:ch33-heap-graph} ... \end{figure}).
% Note: tikzpicture body below is verbatim from the source; do not reformat.
\begin{tikzpicture}[node distance=1.0cm and 1.5cm]
\node[basic block] (r1) {root R1};
\node[basic block, below=of r1] (a) {A};
\node[basic block, below left=of a] (b) {B};
\node[basic block, below right=of a] (c) {C};
\node[basic block, right=of r1] (r2) {root R2};
\node[basic block, below=of r2] (d) {D};
\node[basic block, below=of d] (e) {E dead};
\node[basic block, left=of e] (f) {F dead};
\draw[cfg edge] (r1) -- (a);
\draw[cfg edge] (a) -- (b);
\draw[cfg edge] (a) -- (c);
\draw[cfg edge] (r2) -- (d);
\draw[cfg edge] (d) -- (a);
\end{tikzpicture}
 inside a figure environment
%   (wrap with \begin{figure} ... \caption{...} \label{fig:ch33-heap-graph} ... \end{figure}).
% Note: tikzpicture body below is verbatim from the source; do not reformat.
\begin{tikzpicture}[node distance=1.0cm and 1.5cm]
\node[basic block] (r1) {root R1};
\node[basic block, below=of r1] (a) {A};
\node[basic block, below left=of a] (b) {B};
\node[basic block, below right=of a] (c) {C};
\node[basic block, right=of r1] (r2) {root R2};
\node[basic block, below=of r2] (d) {D};
\node[basic block, below=of d] (e) {E dead};
\node[basic block, left=of e] (f) {F dead};
\draw[cfg edge] (r1) -- (a);
\draw[cfg edge] (a) -- (b);
\draw[cfg edge] (a) -- (c);
\draw[cfg edge] (r2) -- (d);
\draw[cfg edge] (d) -- (a);
\end{tikzpicture}
 inside a figure environment
%   (wrap with \begin{figure} ... \caption{...} \label{fig:ch33-heap-graph} ... \end{figure}).
% Note: tikzpicture body below is verbatim from the source; do not reformat.
\begin{tikzpicture}[node distance=1.0cm and 1.5cm]
\node[basic block] (r1) {root R1};
\node[basic block, below=of r1] (a) {A};
\node[basic block, below left=of a] (b) {B};
\node[basic block, below right=of a] (c) {C};
\node[basic block, right=of r1] (r2) {root R2};
\node[basic block, below=of r2] (d) {D};
\node[basic block, below=of d] (e) {E dead};
\node[basic block, left=of e] (f) {F dead};
\draw[cfg edge] (r1) -- (a);
\draw[cfg edge] (a) -- (b);
\draw[cfg edge] (a) -- (c);
\draw[cfg edge] (r2) -- (d);
\draw[cfg edge] (d) -- (a);
\end{tikzpicture}
